\part{VI: Local Logic and Global Obstruction}\label{part-vi-local-logic-and-global-obstruction}

\hypertarget{chapter-20-local-support-and-logical-regimes}{%
\chapter*{Chapter 20 --- Local Support and Logical Regimes}
\addcontentsline{toc}{chapter}{Chapter 20 --- Local Support and Logical Regimes}
\markboth{Chapter 20 --- Local Support and Logical Regimes}{Chapter 20 --- Local Support and Logical Regimes}\label{chapter-20-local-support-and-logical-regimes}}

A proposition is usually introduced as something that is simply true or
false, with the four-valued and gap-tolerant systems arriving later as
modifications to that starting point --- extra machinery bolted onto a
prior notion of truth value. This chapter takes the opposite route.
Truth values are not primitive here. They are read off something more
basic: what evidence is actually available for and against a claim.

\textbf{Support, not truth value, as the starting object.} Fix a proposition
\(p\) and a region \(U\) of whatever domain is under discussion --- a context,
an observational patch, a tile in the sense already established elsewhere
in this book. Rather than asking whether \(p\) is true at \(U\), ask two
separate questions: is there positive evidence for \(p\) at \(U\), and is
there negative evidence for \(p\) at \(U\)? These questions are independent
of one another --- nothing forces an answer to one to constrain the answer
to the other --- and that independence is the whole of what makes four
values rather than two arise naturally, without stipulation.

Formally, let \(\mathcal{E}^+(U)\) and \(\mathcal{E}^-(U)\) be free abelian
groups generated by the distinct evidential witnesses --- observations,
derivations, testimony --- available for and against \(p\) within \(U\)
respectively. A single patch's support is then the pair
\((e^+, e^-) \in \mathcal{E}^+(U) \times \mathcal{E}^-(U)\), and the
\textbf{support map}

\[q(e^+, e^-) = \begin{cases} N & e^+ = 0,\ e^- = 0 \\ T & e^+ \neq 0,\ e^- = 0 \\ F & e^+ = 0,\ e^- \neq 0 \\ B & e^+ \neq 0,\ e^- \neq 0 \end{cases}\]

reads off one of Belnap and Dunn's~\cite{belnap1977,dunn1976} four values --- Neither, True, False,
Both --- as a consequence of the support data, not as a starting
assumption. \(\mathbb{F} = \{N,T,F,B\}\), ordered by the information order
\(N \sqsubseteq T \sqsubseteq B\), \(N \sqsubseteq F \sqsubseteq B\) with \(T,F\)
incomparable, is a complete lattice, and this matters later; for now it
is enough to note that \(q\) is a genuine function of the underlying
evidence, and that a glut is simply what you get when both kinds of
evidence happen to be present. Nothing about tolerating contradiction has
been stipulated. The four values are a description of four possible
evidential situations, and Both is one of them, on exactly the same
footing as the other three.

\textbf{What this chapter does not yet ask.} Everything above is stated for a
single patch \(U\) in isolation. It says nothing about what happens when
two overlapping patches each report support for the same \(p\) --- whether
their evidence agrees, whether it should be expected to, or what follows
if it doesn't. That is a genuinely separate question, deferred entirely
to Chapter 21. Keeping it separate is not merely tidy exposition; the
distinction between what a single patch supports and what multiple
patches agree on turns out to be the load-bearing distinction of this
part of the book, and blurring it --- treating ``every patch says \(T\)'' and
``every patch says \(T\) for the same reason'' as the same fact --- is
precisely the mistake this book's formal apparatus is built to make
impossible to state by accident.

\begin{center}\rule{0.5\linewidth}{0.5pt}\end{center}

\hypertarget{chapter-21-compatibility-and-gluing}{%
\chapter*{Chapter 21 --- Compatibility and Gluing}
\addcontentsline{toc}{chapter}{Chapter 21 --- Compatibility and Gluing}
\markboth{Chapter 21 --- Compatibility and Gluing}{Chapter 21 --- Compatibility and Gluing}\label{chapter-21-compatibility-and-gluing}}

Chapter 20 fixed what a single patch supports. This chapter asks whether
neighboring patches' support for the same proposition is the \emph{same}
support --- and gets the answer from a much simpler place than it might
appear to require.

\textbf{The question, precisely.} Let \(\{U_i\}\) be a cover, and suppose each
patch \(U_i\) reports local evidence \(e^+_i \in \mathcal{E}^+(U_i)\) for
\(p\)'s positive case (symmetrically for \(e^-_i\)). Where two patches
overlap, restriction maps \(\rho_{U_{ij}, U_i}\) carry each patch's evidence
down to what's visible on the shared region \(U_{ij} = U_i \cap U_j\). The
question is whether \(\rho_{U_{ij},U_i}(e^+_i)\) and \(\rho_{U_{ij},U_j}(e^+_j)\)
agree, as elements of \(\mathcal{E}^+(U_{ij})\).

It is tempting, and was in fact this book's own first attempt, to reach
here for Čech cohomology and ask whether some class is nonzero in
\(H^1\). That reach is wrong, for a reason worth stating once, precisely,
so it need not be revisited: for any specific chosen tuple of local
sections \(s = (s_i)\), the comparison \(\delta s\) --- the collection of
pairwise differences on overlaps --- is by construction a coboundary, and
therefore trivial in cohomology, whether or not the \(s_i\) actually agree.
Asking ``do these particular witnesses match'' is answered directly by
whether \(\delta s\) itself is the zero cochain, not by any cohomology
class built from it. The right tool is simpler than the wrong one, not
more elaborate, and this chapter uses only the tool the question actually
needs.

\textbf{Coherence defect.} Define
\[\delta^+(p) = \big(\rho_{U_{ij},U_j}(e^+_j) - \rho_{U_{ij},U_i}(e^+_i)\big)_{ij}\]
as an element of the Čech 1-cochains of \(\mathcal{E}^+\) relative to the
cover, and \(\delta^-(p)\) symmetrically. Say the positive (respectively
negative) evidence for \(p\) is \textbf{polarity-coherent} across the cover when
\(\delta^+(p) = 0\) (respectively \(\delta^-(p) = 0\)): every pair of
overlapping patches is contributing the \emph{same} witnesses, not merely
witnesses that happen to license the same verdict.

\begin{figure}[htbp]
\centering
\begin{tikzpicture}[
  node distance=1.1cm,
  font=\small
]
\node (title1) {Gluing};
\node[below=0.5cm of title1] (s1) {$s_1$};
\node[right=1.6cm of s1] (s2) {$s_2$};
\node[below=0.9cm of $(s1)!0.5!(s2)$] (ov) {$s_1|_{U_{12}} = s_2|_{U_{12}}$};
\node[below=0.9cm of ov] (glued) {$s \in \mathcal{A}(U_1 \cup U_2)$};
\draw[-{Stealth[length=2mm]}] (s1) -- (ov);
\draw[-{Stealth[length=2mm]}] (s2) -- (ov);
\draw[-{Stealth[length=2mm]},double] (ov) -- (glued);

\begin{scope}[xshift=6.3cm]
\node (title2) {Information-preserving merge};
\node[below=0.9cm of title2] (T) {$T$};
\node[right=1.6cm of T] (F) {$F$};
\node[below=0.9cm of $(T)!0.5!(F)$] (Bv) {$B$};
\draw[-{Stealth[length=2mm]}] (T) -- (Bv);
\draw[-{Stealth[length=2mm]}] (F) -- (Bv);
\end{scope}
\end{tikzpicture}
\caption{Gluing (left) requires a matching family: local sections that
already agree on the overlap combine into a unique global section.
Information-preserving merge (right) requires no such agreement: an FDE
join can combine \emph{disagreeing} reports (positive-only and
negative-only) into $B$, retaining rather than discarding what conflicts.
These are different operations, not two names for one --- see the
retracted correspondence discussed below.}
\label{fig:gluing-vs-merge}
\end{figure}

\textbf{Two different senses of restriction.} The restriction maps used above
carry a modeling choice worth making explicit rather than leaving
implicit, since the two readings genuinely differ. Restriction can mean
plain domain restriction --- the same assignment, viewed on a smaller
region, nothing lost. Or it can mean projection with possible
information loss --- evidence available across the whole of \(U_i\) need not
remain reconstructible from \(V \subset U_i\) alone, if the witnessing
particulars live outside \(V\)'s own resources. The second reading is the
one under which the coherence defect is doing real epistemic work rather
than bookkeeping, and it is the reading this book adopts going forward:
restriction is a genuine act of observation-with-possible-loss, not a
notational convenience. Spatial restriction (the morphism \(\rho_{VU}\))
and epistemic refinement (an ordering of increasing information, as in
Kripke-style forcing~\cite{kripke1963}) are different things, and treating them as
interchangeable was an earlier error in this construction's history worth
naming so it does not recur silently.

\textbf{What this chapter still does not settle.} Coherence, as defined here,
says nothing about which FDE value the patches end up reporting. A
proposition can be perfectly polarity-coherent and still be a glut (every
patch agrees, using the same witnesses, that \(p\) is Both). A proposition
can be badly incoherent and still show no glut at all --- every patch might
independently report \(T\), using witnesses that don't match each other in
the slightest. Whether that combination is even possible, and what it
would mean, is Chapter 22's question, not this one's --- this chapter has
deliberately stopped at the evidence layer and gone no further.

\begin{center}\rule{0.5\linewidth}{0.5pt}\end{center}

\hypertarget{chapter-22-contradiction-is-not-obstruction}{%
\chapter*{Chapter 22 --- Contradiction Is Not Obstruction}
\addcontentsline{toc}{chapter}{Chapter 22 --- Contradiction Is Not Obstruction}
\markboth{Chapter 22 --- Contradiction Is Not Obstruction}{Chapter 22 --- Contradiction Is Not Obstruction}\label{chapter-22-contradiction-is-not-obstruction}}

Chapters 20 and 21 built two structures over the same local data: what a
patch supports (Chapter 20), and whether neighboring patches' support
coheres (Chapter 21). This chapter's claim is that these are not two
views of one fact. They are genuinely separate facts, capable of varying
independently --- and that independence, rather than any single dramatic
case of contradiction, is the central result of this part of the book.

\textbf{The merge.} Given a family of local reports \(\{q_i(p)\}\), define the
pointwise merge
\[\Sigma(p) = \bigvee_i q_i(p),\]
the join in \(\mathbb{F}\)'s information order. Because \(\mathbb{F}\) is a
complete lattice, this join always exists --- there is no family of local
reports the merge fails to combine. Two propositions about it settle what
``merging'' produces:

\begin{quote}
\textbf{Support-Merge Proposition.} \(\Sigma(p) = B\) if and only if some
patch in the family reports positive support and some patch (possibly
the same one) reports negative support. \(\Sigma(p) = N\) if and only if
every patch touching \(p\) reports \(N\).
\end{quote}

\begin{quote}
\textbf{Glut-Free Reconstruction Obstruction.} A pointwise glut-free common
extension of the family into \(\{N,T,F\}\) exists if and only if no
proposition collects both positive and negative support across the
family --- that is, iff \(\Sigma(p) \neq B\) for every \(p\).
\end{quote}

Both are purely combinatorial statements about \(q_i\)-values; neither
mentions the coherence defect from Chapter 21. That is not an oversight ---
it is the point. And a genuine asymmetry falls out of the merge that is
worth stating on its own: \(N\) merges with anything to return that thing
unchanged (\(N \vee x = x\)), so absence of evidence never itself obstructs
reconstruction. Only actually incompatible polar evidence does. Gluts and
gaps are not mirror images of one another, whatever FDE's fourfold
notation suggests at a glance: \textbf{a glut is incompatible polar evidence; a
gap is merely insufficient polar evidence.}

\textbf{Mutual Non-Determination Proposition.} Neither the coherence status
\((\delta^+(p), \delta^-(p))\) from Chapter 21 nor the merged value
\(\Sigma(p)\) from this chapter determines the other. The proof is a
four-cell construction, and its economy is the whole argument: a
two-patch cover with overlap \(U_{12}\) populates all four combinations
directly.

\emph{A note on continuity with Part III.} This is the same construction
Chapter 12 built at the level of distinct evidence sources rather than
cover-indexed patches --- \(M(p) = \bigvee_i s_i(p)\) there, \(\Sigma(p) = \bigvee_i q_i(p)\) here. The two aren't merely analogous; \(\Sigma\) \emph{is}
\(M\), specialized to the case where the sources \(s_i\) happen to be
indexed by a genuine cover with overlaps, which is exactly what licenses
asking the coherence question (\(\delta^\pm\)) that Chapter 12 had no
apparatus for and deliberately didn't attempt. The name change from \(M\)
to \(\Sigma\) marks that shift in structure, not a change in what the
operation does.

\begin{longtable}[]{@{}
  >{\raggedright\arraybackslash}p{(\columnwidth - 4\tabcolsep) * \real{0.3333}}
  >{\raggedright\arraybackslash}p{(\columnwidth - 4\tabcolsep) * \real{0.3333}}
  >{\raggedright\arraybackslash}p{(\columnwidth - 4\tabcolsep) * \real{0.3333}}@{}}
\toprule\noalign{}
\begin{minipage}[b]{\linewidth}\raggedright
\end{minipage} & \begin{minipage}[b]{\linewidth}\raggedright
\(\delta\) coherent
\end{minipage} & \begin{minipage}[b]{\linewidth}\raggedright
\(\delta\) incoherent
\end{minipage} \\
\midrule\noalign{}
\endhead
\bottomrule\noalign{}
\endlastfoot
\textbf{\(\Sigma(p) = T\)} & (a) both patches report \(T\) via the \emph{same} witness & (c) both patches report \(T\) via \emph{different}, non-matching witnesses \\
\textbf{\(\Sigma(p) = F\)} & (a\('\)) both patches report \(F\) via the \emph{same} witness & (c\('\)) both patches report \(F\) via \emph{different}, non-matching witnesses \\
\textbf{\(\Sigma(p) = B\)} & (b) both polarities individually coherent, genuinely in conflict & (d) conflicting \emph{and} internally incoherent \\
\end{longtable}

The \(T\)/\(F\) rows are genuinely separate cells, not mirror images assumed
free --- (a\('\)) sets \(e^-_1 = e^-_2 = v\) (a shared negative witness, no
positive evidence anywhere), giving \(\delta^- = 0\); (c\('\)) sets
\(e^-_1 = v_1 \neq v_2 = e^-_2\), giving \(\delta^- \neq 0\) while both patches
still report the same verdict \(F\). Both are built the same way (a) and
(c) were, not inferred from them. The \(B\) row already covers both
polarities by construction --- a glut requires both positive and negative
support present somewhere in the family, so (b) and (d) don't need
separate positive- and negative-side versions the way the non-glut rows
did. Six cells, all built, none left to symmetry. That every cell is
populated is the entire proof --- neither axis constrains the other,
because both are demonstrably free to vary while the other is held fixed.

\begin{figure}[htbp]
\centering
\begin{tikzpicture}[
  font=\small
]
\matrix (grid) [
  matrix of nodes,
  nodes={draw, minimum width=3.6cm, minimum height=1.3cm, align=center, anchor=center},
  row sep=-\pgflinewidth, column sep=-\pgflinewidth
] {
\shortstack{(a) $\Sigma=T$\\coherent} & \shortstack{(c) $\Sigma=T$\\incoherent} \\
\shortstack{(a$'$) $\Sigma=F$\\coherent} & \shortstack{(c$'$) $\Sigma=F$\\incoherent} \\
\shortstack{(b) $\Sigma=B$\\coherent} & \shortstack{(d) $\Sigma=B$\\incoherent} \\
};
\node[font=\small\itshape, above=0.2cm of grid-1-1.north, anchor=south] {$\delta$ coherent};
\node[font=\small\itshape, above=0.2cm of grid-1-2.north, anchor=south] {$\delta$ incoherent};
\node[font=\small\itshape, rotate=90, anchor=south] at ($(grid.west) + (-0.6,0)$) {merged value $\Sigma(p)$};
\end{tikzpicture}
\caption{The six realizability witnesses proving the Mutual
Non-Determination Proposition. Reading across a row: the same merged
value $\Sigma(p)$ occurs under both coherent and incoherent evidence.
Reading down a column: the same coherence status occurs under all three
merged values. Neither axis is a function of the other. Case (c) is the
one Priest's semantics alone cannot see: a classical-looking verdict
resting on evidence that does not actually agree with itself.}
\label{fig:mutual-non-determination}
\end{figure}

Case (b) is Priest's glut in its cleanest form: two well-founded,
internally coherent positions in genuine disagreement --- the positive case
coheres into one witness, the negative case coheres into a different one,
and they simply conflict. Call this the \textbf{clean glut}. Case (d) is the
same conflict compounded by internal incoherence on at least one side ---
part of what looks like cross-polarity disagreement may really be several
non-agreeing positive (or negative) sub-cases masquerading as one. Call
this the \textbf{confounded glut}. Both are legitimate uses of the word
``glut,'' and this book will keep them distinguished rather than treating
every \(\Sigma(p)=B\) as the same kind of event.

Case (c) deserves the most attention of the four, precisely because it is
the least dramatic. Nothing here looks unusual under FDE: both patches
say \(T\), no contradiction anywhere, nothing for a paraconsistent logic to
even notice. And yet the agreement is hollow in a specific, checkable
sense --- the patches are not agreeing about the same thing, only about the
same verdict. This is the case that carries the rest of Part VI, and it
recurs through Chapters 21 to 23 for that reason: it is the demonstration
that the support layer built in Chapters 20 and 21 is not decorative
machinery sitting underneath an otherwise complete four-valued logic. It
is detecting something four-valued logic alone cannot see at all.

\textbf{The headline result.} Contradiction, in Priest's sense, is a fact
about \(\Sigma(p)\). Obstruction to coherent reconstruction, in the sense
this book has been building since Chapter 21, is a fact about
\(\delta^\pm(p)\). This chapter's title is meant literally:
\textbf{contradiction is not obstruction.} A glut can be perfectly coherent. A
non-glut can be badly incoherent. Priest was right that gluts needn't
explode; this chapter's addition is that whether something needs
repairing at all cannot be read off \(\Sigma(p)\) in isolation --- which sets
up exactly the question Chapter 23 has to answer.

\begin{center}\rule{0.5\linewidth}{0.5pt}\end{center}

\hypertarget{chapter-23-repair-and-reconstruction}{%
\chapter*{Chapter 23 --- Repair and Reconstruction}
\addcontentsline{toc}{chapter}{Chapter 23 --- Repair and Reconstruction}
\markboth{Chapter 23 --- Repair and Reconstruction}{Chapter 23 --- Repair and Reconstruction}\label{chapter-23-repair-and-reconstruction}}

The progression through this part of the book has been

\[\text{support (Ch. 19)} \to \text{coherence (Ch. 20)} \to \text{non-determination (Ch. 21)} \to \text{repair (this chapter)}.\]

Chapter 22 established that support and coherence are independent facts.
This chapter asks what a repair operator does with both of them at once
--- and the answer is that it needs a third term it didn't previously have.

\textbf{Why case (c) cannot be folded into existing repair cost.} This book's
repair theory already has a cost formulation, \(E(a) + \lambda R(a)\),
where \(E\) is the ordinary cost of an action \(a\) and \(R\) is the cost of
repairing a polarity conflict --- a glut. Case (c) from Chapter 22 has
\(R(a) = 0\), exactly and correctly: there is no glut, no polarity conflict
of any kind, both patches agree \(p\) is \(T\). Whatever case (c) costs, it
is not a repair cost in the sense \(R\) was built to capture, and treating
it as one would erase the very distinction Chapters 21 and 22 exist to
establish.

\textbf{The third term.} Repair cost needs

\[J(a) = E(a) + \lambda R(a) + \mu C(a),\]

where \(C(a,p)\) is \textbf{coherence liability} --- the cost an action incurs
from depending on evidence that failed to cohere at Chapter 21's level,
even when Chapter 22's merge saw nothing wrong. Crucially, \(C\) cannot be
a fixed function of the coherence defect alone:

\[C(a,p) = \rho_a\big(\delta^+(p), \delta^-(p)\big),\]

action-relative by construction. The coherence defect is an objective
structural fact about the evidence configuration, unchanged by anyone's
plans; \(\rho_a\) is where the actual liability lives, and it depends on
what the action in question needs from the evidence.

A bridge reported open by camera evidence in one context and by
traffic-authority report in another is a witness mismatch --- but for an
action like \emph{drive across}, that mismatch is corroboration, not
incoherence: two independent reasons for the same conclusion. The same
structural fact, applied to an action like \emph{use this credential} where
the two contexts' positive evidence for ``the key is valid'' trace to two
different underlying keys, is disqualifying --- the action needs the
witnesses to be identical, not merely to agree on a verdict. Same
\(\delta^+(p) \neq 0\), opposite consequence, depending entirely on what the
action depends on. This is why \(\rho_a\) cannot be dispensed with in favor
of a single fixed penalty function: the liability is a fact about the
action's relationship to the evidence, not a fact about the evidence
alone.

\textbf{What this buys.} Under the old two-term functional, an action relying
on case (c)'s proposition is indistinguishable from one relying on case
(a) whenever their ordinary costs are equal --- neither has a glut, so
\(R = 0\) for both, and \(E + \lambda R\) has nothing left to compare. Under
the three-term functional, a witness-sensitive action pays through \(C\)
where a witness-insensitive one does not, even though \(\Sigma(p) = T\)
identically for both:

\begin{quote}
Two states indistinguishable under the FDE polarity merge can induce
different action preferences once evidential coherence enters the
decision functional.
\end{quote}

Which gives this part of the book its actual thesis, stated plainly:
\textbf{logical agreement does not imply evidential coherence, and rational
action may depend on both independently.} It is case (c), not the
dramatic glut cases, that earns the extra machinery its place --- a
classical-looking, contradiction-free agreement can still be the wrong
thing to act on, for reasons a four-valued semantics alone has no way to
express.

\textbf{A deliberate incompleteness.} This chapter does not commit to a
canonical form for \(\rho_a\) beyond the minimal binary flag --- sensitive or
insensitive --- used in the bridge/credential contrast above. That is a
choice, not a gap: establishing that an action-relative coherence
liability distinguishes configurations invisible to the old functional is
the result this chapter needs, and it holds regardless of which
particular metric or coefficient structure eventually fills in \(\rho_a\).
Committing to a graded, numerical version now would answer a question
this chapter hasn't yet earned the right to ask --- whether \(\rho_a\) should
depend on \emph{which} witness-identity condition an action needs, and how \(C\)
should aggregate across actions depending on several propositions at
once, are left open for exactly that reason.

\begin{center}\rule{0.5\linewidth}{0.5pt}\end{center}

\hypertarget{status-of-this-draft}{%
\section{Status of this draft}\label{status-of-this-draft}}

Written sequentially, as a test of whether the five revisions of scratch
work compose without the correction history as scaffolding. They do. The
mirror-case gap flagged in the first pass --- \(\Sigma(p) = F\) under each
coherence status, previously asserted by symmetry rather than built --- is
now closed with explicit constructions (a\('\)) and (c\('\)) added to Chapter
21's table. The progression from support to coherence to non-determination
to repair holds together as stated, case (c) does the work across three
chapters that it was meant to, and Chapter 23's deliberate incompleteness
around \(\rho_a\) is the honest state of that question rather than a
placeholder for something already decided.
